% Compile with: uplatex file.tex && dvipdfmx file.dvi
%           or: latexmk -uplatex file.tex
\documentclass[dvipdfmx]{article}
\usepackage{amsmath, amssymb}
\usepackage{tikz}

\title{scenario\_ac\_impedance}
\date{\today}

\begin{document}
\maketitle

To work with a series RLC circuit's impedance, prepare the resistance R, inductance L, angular frequency w (omega), and capacitance C as real scalar parameters.

Define R, L, w, C as scalars

\[
  R, L, w, C
\]

The impedance of a resistor R, an inductor L, and a capacitor C connected in series is the complex quantity Z = R + j*omega*L + 1/(j*omega*C), where j is the imaginary unit (in this electrical\_engineering domain, following engineering convention, j denotes the imaginary unit and i denotes current).

$Z = R + j w L + 1/(j w C)$

\begin{equation}
  Z = R + j w L + 1/(j w C)
  \label{eq:E1}
\end{equation}

Extract the real part of the impedance Z (the resistive component).

Find the realpart of E1

\begin{equation}
  R
  \label{eq:E2}
\end{equation}

Extract the imaginary part of the impedance Z (the reactive component), representing the phase shift between current and voltage caused by the energy-storing elements L and C.

Find the imagpart of E1

\begin{equation}
  L w-\frac{1}{C w}
  \label{eq:E3}
\end{equation}

Resonance is the condition where the reactance (the imaginary part) vanishes. Form the equation that sets this reactance to zero.

$E3 = 0$

\begin{equation}
  L w-1/(C w) = 0
  \label{eq:E4}
\end{equation}

Solving this equation for the angular frequency w gives the resonant angular frequency. Of the two roots that differ in sign, the physically meaningful one is the positive root omega\_0 = 1/sqrt(LC).

Solve E4 for w

\begin{equation}
  w = -\sqrt{\frac{1}{C L}}
  \label{eq:E5}
\end{equation}

\begin{equation}
  w = \sqrt{\frac{1}{C L}}
  \label{eq:E6}
\end{equation}

\end{document}